\section{Тестирование}
\label{sec:tests}

Автотесты написаны на \texttt{unittest} из стандартной библиотеки и лежат в
каталоге \texttt{tests/}. Все они работают без интернета: HTTP-клиент подменяется
классом \texttt{FakeFetcher}, который отдаёт страницы из словаря и, как
настоящий сервер, отвечает \texttt{304}, если \texttt{ETag} совпал. Для разбора
используются настоящие страницы из образцов первой лабораторной работы, а фильтры
проверяются на рабочем \texttt{config.yaml}.

База запускается в двух режимах. \texttt{make test} --- MongoDB в памяти
(\texttt{mongomock}): не нужны ни сеть, ни докер, тесты идут за десятые доли
секунды. \texttt{make test-mongo} --- те же тесты на настоящей MongoDB из
докера: каждый тест получает свою временную базу со случайным именем и удаляет
её после себя. Второй режим нужен, потому что эмулятор повторяет MongoDB не
полностью --- например, он падал на пакетной вставке через \texttt{bulk\_write},
которая на настоящем сервере работала. Приложенный результат получен на
настоящей MongoDB.

\subsection*{План тестирования}

План построен по требованиям задания и по механизмам робота: на каждое
требование --- хотя бы один тест, который упадёт, если требование нарушено.

\begin{center}
\begin{longtable}{|>{\raggedright\arraybackslash}p{4.8cm}|>{\raggedright\arraybackslash}p{8.2cm}|r|}
\hline
\textbf{Что проверяется} & \textbf{Как} & \textbf{Тестов} \\
\hline
\endhead
Нормализация URL (\texttt{test\_urls})
  & склейка относительных ссылок, отбрасывание якоря и служебных параметров,
    раскрытие процентных последовательностей, но не \texttt{\%25} и разделителей,
    регистр, порт, повторные слэши, отказ от \texttt{mailto:}/\texttt{javascript:},
    идемпотентность & 10 \\
\hline
Кодировка (\texttt{test\_parsing})
  & приоритет \texttt{Content-Type} над \texttt{<meta>}, \texttt{<meta charset>} и
    \texttt{http-equiv}, откат на UTF-8, неизвестная кодировка, настоящая страница
    в windows-1251 & 6 \\
\hline
Выделение текста (\texttt{test\_parsing})
  & баланс вложенных \texttt{div}, порядок атрибутов, выбор контейнера по
    \texttt{itemprop}, удаление скриптов и таблиц, \texttt{cut\_after}; на
    настоящих страницах Википедии, film.ru и Кино-Театр.Ру --- отсутствие меню,
    подвала, подписи автора и кнопок; раскрытие HTML-сущностей & 12 \\
\hline
Ссылки, заголовок, отпечаток (\texttt{test\_parsing})
  & нормализация и дедупликация ссылок, HTML-сущности в \texttt{href},
    ссылки настоящей страницы, заголовок, SHA-256 & 5 \\
\hline
Фильтры (\texttt{test\_rules})
  & на рабочем конфиге. Википедия: обходятся только категории о кино,
    мета-категории и адреса с параметрами (правка, красные ссылки,
    \texttt{/w/}) отсекаются, категории не сохраняются, со статей ссылки не
    собираются. film.ru: статьи сохраняются, страницы списков обходятся, прочее
    отсекается, на настоящей странице списка находятся все 28 статей. Признак
    капчи & 13 \\
\hline
Хранилище (\texttt{test\_storage})
  & уникальность в очереди, порядок выдачи, выбор источника, приоритет
    переобкачки важнее выбора источника, пакетная вставка, уникальный индекс, ревизии, обязательные поля
    документа, выбор документов для проверки без тел, очередь переживает
    переподключение к базе; переразбор текста из \texttt{raw\_html} без
    перекачки & 14 \\
\hline
Робот целиком (\texttt{test\_crawler})
  & полный обход: сохраняется только документ, служебная и короткая страницы
    отсеяны; в базе только текст контейнера; \texttt{max\_depth},
    \texttt{max\_documents}; \texttt{robots.txt}; ошибка 404; \textbf{обрыв
    сети --- повтор с паузой}, долгий обрыв --- \texttt{failed} после пяти
    попыток; капча --- пауза источника с нарастанием и повтор; шаблон стартовых
    страниц; \textbf{остановка и продолжение с того же URL}; повторный запуск
    без дублей & 13 \\
\hline
Выбор источника (\texttt{test\_crawler})
  & хост с неистёкшей паузой не ждём, а берём другой источник; готовые
    источники чередуются & 2 \\
\hline
Переобкачка целиком (\texttt{test\_crawler})
  & \texttt{304} по \texttt{ETag}; изменилась разметка, но не текст --- документ
    не трогаем; изменился текст --- URL в очереди с приоритетом, робот
    перекачивает его первым и поднимает \texttt{revision}; недавно
    проверенные документы не запрашиваются & 4 \\
\hline
\multicolumn{2}{|r|}{\textbf{Всего}} & \textbf{79} \\
\hline
\end{longtable}
\end{center}

Остановка проверяется так: флаг остановки срабатывает после первого URL, затем
к той же базе подключается новый экземпляр робота. Тест требует,
чтобы второй запуск запросил ровно оставшиеся адреса и не запросил повторно
уже обработанный.

\subsection*{Результат}

\begin{verbatim}
$ make test-mongo
...
test_stop_and_resume (tests.test_crawler.CrawlerTest....) ... ok
...
test_queue_survives_reopen (tests.test_storage.PersistenceTest...) ... ok
----------------------------------------------------------------------
Ran 79 tests in 3.351s

OK
\end{verbatim}

Полный вывод --- в файле \texttt{report/tests.log} (\texttt{make test-report}).

\subsection*{Проверка базы}

Тесты проверяют код, но не то, что лежит в базе после настоящего прогона. Для
этого есть отдельная утилита, \texttt{make validate}. По каждому документу она
проверяет:

\begin{itemize}
    \item все обязательные поля на месте, \texttt{crawled\_at} --- целое число
          секунд в разумных пределах, источник есть в конфиге;
    \item адрес уникален, нормализован (повторная нормализация его не меняет)
          и проходит фильтры своего источника;
    \item выделенный текст заново получается из сохранённого \texttt{raw\_html}
          символ в символ, а \texttt{content\_hash} равен SHA-256 текста, ---
          значит, текст взят именно из этой страницы;
    \item текст не короче порога, капчи среди документов нет;
    \item каждый документ отмечен в очереди как обработанный.
\end{itemize}

С ключом \texttt{LIVE=N} утилита ещё перекачивает $N$ случайных документов с
сайтов и сравнивает с базой --- так видно, что документы настоящие и доступны.
На итоговом корпусе все проверки пройдены:

\begin{verbatim}
Документов: 79566, очередь: {'skipped': 554, 'failed': 6388, 'done': 125228}
[ok  ] все обязательные поля на месте: 79566/79566
[ok  ] url уникальны: 79566/79566
[ok  ] url нормализованы: 79566/79566
[ok  ] url проходит фильтры источника: 79566/79566
[ok  ] text заново выделяется из raw_html один в один: 79566/79566
[ok  ] content_hash = SHA-256 текста: 79566/79566
[ok  ] нет страниц-капч: 79566/79566
...
Итог: всё в порядке
\end{verbatim}

Полная проверка 79 тысяч документов занимает несколько минут: утилита читает
базу потоком, по одному документу, и не держит её в памяти целиком.

\subsection*{Что нашли тесты}

Тесты на настоящих страницах нашли две ошибки в выделении текста и ссылок,
которые прогон робота не показал. Третью показал прогон.

\textbf{Контейнер Кино-Театр.Ру не находился.} Регулярное выражение искало тег
ровно в виде \texttt{<div class="...">}, а на сайте он записан как
\texttt{<div itemprop="text" class="art article\_text">}. Контейнер не
находился, и в текст попадала вся страница с меню и подвалом. Вдобавок даже
найденный контейнер брался до конца страницы, а не до парного
\texttt{</div>}. У биографий и аннотаций контейнера \texttt{art article\_text}
нет вовсе, текст там лежит в \texttt{<div itemprop="description">}.

После исправления текст статьи о \enquote{Месте встречи} сократился с 21\,280
до 12\,852 символов, а биографий актёров --- с 8--9 тысяч до 200--400
символов. Отсюда неприятный вывод: 17 документов Кино-Театр.Ру из первого прогона
учтены вместе с навигацией, и средний объём текста 21\,240 символов по этому
источнику завышен. Биографии прошли порог \texttt{min\_text\_length} только за счёт
меню. С исправленным выделением они в корпус не попадают: биографии на этом
сайте слишком короткие, чтобы быть документами. Это вторая, после блокировки,
причина, по которой источник в итоге заменён на film.ru.

\textbf{Знак процента в адресе.} Уже демонстрационный прогон на настоящих сайтах
показал адрес \texttt{/wiki/1\_\%\_(фильм)}. Нормализация раскрывала все
процентные последовательности, включая \texttt{\%25} --- экранированный сам знак
процента, --- и адрес становился неоднозначным: голый \texttt{\%} в пути
некорректен. Теперь \texttt{\%25}, \texttt{\%2F}, \texttt{\%3F} и \texttt{\%23}
не раскрываются, на это добавлен тест.

\textbf{HTML-сущности в ссылках.} В атрибуте \texttt{href} амперсанд записан как
\texttt{\&amp;}, и без раскрытия сущностей адрес
\texttt{index.php?title=A\&amp;action=edit} превращался в параметр
\texttt{amp;action}. На Википедии такие ссылки всё равно отсекаются фильтром, но
на сайтах с параметрами в адресах это давало бы дубли одного документа.

\textbf{HTML-сущности в тексте.} Первая версия раскрывала только \texttt{\&nbsp;},
а остальные сущности удаляла: \texttt{\&laquo;Сталкер\&raquo; \&mdash; \ldots{}
\&amp; Стругацких} превращалось в \enquote{Сталкер \ldots{} Стругацких} --- без
кавычек, тире и амперсанда. Сильнее всего это задело film.ru, который пишет
кавычки сущностями. Теперь сущности раскрываются, а весь корпус переразобран
утилитой \texttt{make reparse} из сохранённого HTML, без единого запроса к сайтам.

Когда правила выделения текста меняются, уже скачанные документы не нужно
качать заново: в базе хранится \enquote{сырой} HTML, и текст переразбирается из
него. На этом же построена проверка \texttt{make validate}: она заново выделяет
текст из \texttt{raw\_html} и сравнивает с сохранённым.

\pagebreak
